\scanpage{36}
% L1–5
$M$ module de type fini sur l'anneau local $A$ de dim $n$

$k$ corps résiduel de $A$

$\unsure{\mathfrak n|\mathfrak N}$ idéal de définition fixé

$x_1\ldots x_n$ syst. de paramètres fixé

$\mathfrak p$ un idéal premier de $A$ de corang $n$, fixé \qquad $N$ variable de long. finie, $N_0\ne0$ fixé

% L6–22, two columns connected by implication arrows
\deleted{\unsure{(iii)|(iv)}}
\diagram{\begin{tikzcd}[ampersand replacement=\&,row sep=.75em,column sep=.3em,font=\small]
\text{(i)}\arrow[d,Rightarrow]\arrow[dddd,bend right=40,Rightarrow] \&
\parbox[t]{12em}{\raggedright $\operatorname{Ext}_A^{n-1}(k,M)=0$ et l'une des conditions \inserted{a) à d)} suivantes [qui sont équivalentes moyennant la première, car on aura $\operatorname{Ext}^{n}(N,M)\simeq\operatorname{Hom}(N,H^n(M))$ etc ]} \& \parbox[t]{18em}{\raggedright a) $\operatorname{Ext}^{n+1}(k,M)=0\qquad \operatorname{Ext}^{n}(k,M)\simeq k$\newline
\deleted{\unsure{b)}} $*$b) $N\to\operatorname{Ext}^{n}(N,M)$ est dualisant\newline
$*$c) $N\to\operatorname{Ext}^{n}(N,M)$ est exact, $0<\lg_{\mathfrak p}M\leq\lg_{\mathfrak p}A$\newline
b') $H^n(M)$ dualisant\newline
c') $H^n(M)$ injectif, $0<\lg_{\mathfrak p}M\leq\lg_{\mathfrak p}A$\newline
d) $\operatorname{Ext}^{n+1}(k,M)=0$, $0<\lg_{\mathfrak p}M\leq\lg_{\mathfrak p}A$}\\
\text{(i bis)}\arrow[u,dotted,Rightarrow,"{\substack{\text{[si }M\,\deleted{\unsure{\text{Coh. Mac.}}}\\\text{\inserted{équidim.}}\,\text{]}}}" description]\arrow[d,Rightarrow] \&
\parbox[t]{12em}{\raggedright $\operatorname{Ext}_{A}^{n-1}(k,M)=0$ et l'une des \deleted{deux cond.} équiv. suivantes} \& \parbox[t]{18em}{\raggedright 
$*$e) $\operatorname{Ext}^{n}(k,M)\simeq k$, $M$ fidèle\newline
$*$f) $\operatorname{Ext}^{n}(A/\mathfrak n,M)$ est comonogène, $M$ fidèle}\\
\text{(ii)}\arrow[d,Rightarrow]\&
\parbox[t]{12em}{\raggedright Sous l'une des formes suivantes, qui sont équivalentes} \& \parbox[t]{18em}{\raggedright 
a) $N\to\operatorname{Ext}^{n}(N,M)$ est dualisant\newline
b) $N\to\operatorname{Ext}^{n}(N,M)$ exact, $0<\lg_{\mathfrak p}M\leq\lg_{\mathfrak p}A$}\\
\text{(iii)}\& \operatorname{Ext}^n(k,M)\simeq k,\quad M\ \text{fidèle}\\
\text{(iv)}\& \operatorname{Ext}^n(A/\mathfrak n,M)\simeq D(A/\mathfrak n),\quad M\ \text{fidèle}
\end{tikzcd}}
% L18–22, right of (iii)–(iv)
\margin{\deleted{\illegible{3}} NB, \unsure{même|Ici} si \inserted{$A$ et} $M$ sont Coh. Mac. il ne suffit pas ici de \unsure{considérer} \deleted{\illegible{4}} la \unsure{longueur}, car $\operatorname{Ext}^n(\unsure{M|N},M)=\operatorname{Hom}(\unsure{M|N},H^n(M))$ en dim $0$ \unsure{p.ex.}, si on ajoute des \unsure{élémts} nilpotents à $A$ \deleted{\illegible{2}}}

% L23–24
peut-être suffit-il même que $\operatorname{Ext}^n(A/\mathfrak n,M)$ soit comonogène et $M$ de support $\operatorname{Spec}(A)$

% L25–26
(v) L'annulateur de $\mathfrak m$ dans $M/x_1M+\ldots+x_nM$ est de dim $1$, $M$ est fidèle

% L27–35
Ces conditions sont vraies si $\boxed{A\text{ est Coh Mac et }M\simeq\Omega_A}$

Pour prouver la réciproque, il suffirait de montrer que dans les cas \unsure{iv)|ii)|(i bis)} à iii) (si on peut \unsure{iii), i)bis!?|(i bis) faible}) que \inserted{si $n>0$,} l'annulateur de $\mathfrak m$ dans $M$ est nul, \deleted{dans le cas (iv), qu'il en est de même dans \unsure{aussi|avec} $M$ et $A$,} Dans (v), \deleted{\unsure{que} si $n>0$ et si} \inserted{peut songer à} commencer par prouver que si $n>0$, et la \uline{première} condition de (v) est vérifiée, alors l'annulateur de $\mathfrak m$ dans $M$ est nul [d'où déjà que $M$ est Coh-Mac par récurrence, puis que $\operatorname{Ext}^n(A/\underline{x},M)\simeq M/\underline{x}M$ et $\operatorname{Ext}^n(k,M)\simeq k$, d'où le fait que $H^n(M)$ est le dual d'un $A/J$, $J$ intersection d'idéaux premiers de corang $n$

% archive sheet number: 35
